% Part: propositional-logic
% Chapter: propositional-logic 
% Section: completeness

\documentclass[../../../include/open-logic-section]{subfiles}

\begin{document}

\olfileid{pl}{prp}{com}

\olsection{Kajangkepan Logika Proposisional}

\begin{defn}
\ollabel{def:MaxCon}
Himpunan $\Gamma$ saka !!{formula}s diarani \emph{konsisten maksimal}
yen konsisten lan yen $\Delta$ minangka himpunan konsisten kanthi
$\Gamma \subseteq \Delta$, mula $\Gamma = \Delta$.
\end{defn}

\begin{lem}[Lema Kabeneran]
  \ollabel{lem:truth}
Upamane $\Gamma$ konsisten maksimal; mula:
\begin{enumerate}
\item \ollabel{lem:truth-1} $!A \in \Gamma$ yen lan mung yen
  $\lnot !A \notin \Gamma$;
\item \ollabel{lem:truth-2} $\Gamma \Proves !A$ yen lan mung yen $!A \in \Gamma$;
\item \ollabel{lem:truth-3} $!A \lif !B \in \Gamma$ yen lan mung yen
  $!A \notin \Gamma$ utawa $!B \in \Gamma$.
\end{enumerate}
\end{lem}

\begin{proof}
\begin{enumerate}
\item Upamane $!A \in \Gamma$. Yen uga $\lnot !A \in \Gamma$,
  mula $\Gamma$ ora konsisten; lan yen $!A$ lan $\lnot!A$
  kalorone ora ana ing $\Gamma$, mula miturut \olref[axd]{prop:phi},
  salah siji saka $\Gamma\cup\{!A\}$ lan $\Gamma \cup \{\lnot !A\}$
  konsisten. Iki ateges $\Gamma$ ora maksimal.
  Arah kosok baline dibuktekake kanthi cara kang padha.
\item Yen $!A \in \Gamma$, mesthi $\Gamma \Proves !A$.
  Kosok baline, upamane $\Gamma \Proves !A$. Yen $!A \notin \Gamma$,
  mula miturut \olref{lem:truth-1}, $\lnot !A \in \Gamma$.
  Dadi $\Gamma \Proves !A$ lan $\Gamma \Proves \lnot !A$, sawijining kontradiksi.
\item Latihan.
\end{enumerate}
\end{proof}

Arah saka kiwa menyang tengen ing \olref{lem:truth}\olref{lem:truth-2}
nuduhake yen $\Gamma$ \emph{katutup tumrap deduksi}, yaiku yen
$\Gamma \Proves !A$, mula $!A \in \Gamma$ kanggo kabeh~$!A$.

\begin{prob}
  Jangkepana bukti \olref{lem:truth}.
\end{prob}

\begin{thm}[Kajangkepan]
\ollabel{thm:completeness}
Yen $\Gamma$ konsisten, mula $\Gamma$ bisa disembadani.
\end{thm}

\begin{proof}
Upamane $!A_0$, $!A_1$,~\dots minangka dhaptar kang nyakup kabeh
!!{formula}s ing basa kasebut. Kita nemtokake kanthi rekursif sawijining
urutan himpunan !!{formula}s $\Gamma_0$, $\Gamma_1$,~\dots kang saya
gedhe, kanthi netepake:
\begin{align*}
\Gamma_0 & = \Gamma\\
\Gamma_{n+1} & =
\begin{cases}
  \Gamma_n \cup \{!A_n\} & \text{yen $\Gamma_n \cup \{!A_n\}$
    konsisten;}\\
  \Gamma_n \cup \{\lnot !A_n\} & \text{yen ora}.
\end{cases}
\end{align*}
Banjur tetepna:
\[
\Gamma^* = \bigcup_{0\le n}\Gamma_n.
\]
Bukti diterusake kanthi netepake kasunyatan ing ngisor iki kanthi urut:
\begin{enumerate}
\item Kanggo saben $n$, himpunan $\Gamma_n$ konsisten
  (kanthi induksi marang $n$, nggunakake \olref[axd]{prop:phi});
\item $\Gamma^*$ konsisten;
\item $\Gamma^*$ maksimal.
\end{enumerate}
Banjur tetepna valuasi $v$ kanthi menehi $v(p_i) = \True$ yen lan
mung yen $p_i \in \Gamma^*$. Kanthi induksi marang $!A$, dituduhake
yen kaanggotaan ing $\Gamma^*$ cocog karo kabeneran miturut $v$
uga kanggo !!{sentence}s kang luwih rumit:
$\overline{v}(!A) = \True$ yen lan mung yen $!A \in \Gamma^*$.
Mligine, $v$ nyembadani $\Gamma^*$. Amarga $\Gamma \subseteq \Gamma^*$,
$\Gamma$ uga bisa disembadani, kaya kang dikarepake.
\end{proof}

\begin{cor}
Yen $\Gamma \Entails !A$, mula $\Gamma \Proves !A$.
\end{cor}

\begin{proof}
Yen $\Gamma\Proves/ !A$, mula $\Gamma \cup \{\lnot !A\}$ konsisten,
miturut \olref[axd]{prop:prov-incons}. Dadi, miturut teorema kajangkepan,
$\Gamma \cup \{\lnot !A\}$ bisa disembadani, lan $\Gamma \Entails/ !A$.
\end{proof}

\begin{prop}[Teorema Kekompakan]
\ollabel{prop:compactness}
$\Gamma$ bisa disembadani yen lan mung yen saben subhimpunan
\emph{winates} $\Gamma_0$ saka $\Gamma$ bisa disembadani.
\end{prop}

\begin{proof}
$\Gamma$ ora bisa disembadani yen lan mung yen ora konsisten,
yen lan mung yen ana subhimpunan winates $\Gamma_0$ saka $\Gamma$
kang ora konsisten, yen lan mung yen ana subhimpunan winates
$\Gamma_0$ saka $\Gamma$ kang ora bisa disembadani.
\end{proof}

\begin{cor}
$\Gamma \Entails !A$ yen lan mung yen ana subhimpunan winates
$\Gamma_0$ saka $\Gamma$ kanthi $\Gamma_0 \Entails !A$.
\end{cor}

\end{document}
